\chapter{Algorithmes de référence}\label{ch:mx:benchmark-algorithms}

La plupart des ordres institutionnels sur actions sont confiés à des algorithmes qui promettent un indice de référence~: le prix moyen pondéré par les volumes de la journée, le prix d'arrivée, le cours de clôture.
Chaque promesse repose sur un échéancier, et chaque échéancier ne vaut que par la prévision de volume de la journée sur laquelle il s'appuie. Ce chapitre construit les cinq algorithmes classiques sur
l'interface d'ordonnanceur du chapitre~14, prévoit le volume intrajournalier de trois façons et fait tourner l'ensemble sur le marché simulé, afin de mesurer dans quelle part de l'écart
d'un algorithme par rapport à son indice de référence est imputable à la prévision de volume, ce que fait un \omterm{def:mx:benchmark-algorithms:vwap}{algorithme de participation} au volume qu'il suit, et comment
un indice de référence peut être déplacé par l'ordre mesuré par rapport à lui.

\section{Les algorithmes et leurs indices de référence}

\begin{definition}[Algorithme d'exécution]\label{def:mx:benchmark-algorithms:algo}
Un \emph{algorithme d'exécution}\index{algorithme d'exécution} est un programme qui traite un \omterm{def:mx:the-almgren-chriss-framework:parent}{ordre parent} pour le compte d'un client ou d'un desk~: il décide de la quantité à négocier
à chaque intervalle (l'échéancier) et envoie les \omterm{def:mx:the-almgren-chriss-framework:parent}{ordres enfants}, mesurés par rapport à un prix de référence fixé à l'avance.
\end{definition}

\begin{definition}[Algorithmes VWAP, TWAP et de participation]\label{def:mx:benchmark-algorithms:vwap}
Un \emph{algorithme VWAP}\index{algorithme VWAP} négocie la quantité de chaque intervalle en proportion de la part du volume de la journée qu'il y attend, visant
le prix moyen pondéré par les volumes de la journée. Un \emph{algorithme TWAP}\index{algorithme TWAP} négocie des quantités égales dans des intervalles égaux, visant le
prix moyen pondéré par le temps. Un \emph{algorithme de participation}\index{algorithme de participation} négocie un pourcentage fixe du volume au fur et à mesure que celui-ci s'affiche,
de sorte que son échéancier suit le marché au lieu de suivre une prévision.
\end{definition}

\begin{definition}[Algorithmes d'implementation shortfall et de clôture cible]\label{def:mx:benchmark-algorithms:is}
Un \emph{algorithme d'implementation shortfall}\index{algorithme d'implementation shortfall} négocie selon un échéancier d'exécution optimale (la trajectoire d'Almgren--Chriss
du chapitre~14, avancée par l'aversion au risque), visant le prix d'arrivée. Un
\emph{algorithme de clôture cible}\index{algorithme de clôture cible} vise le cours de clôture~: il négocie une partie de l'ordre dans l'enchère de clôture au moyen d'un \omterm{def:mx:the-limit-order-book:orders}{ordre au marché} à la clôture (MOC), et le reste dans la dernière
partie de la journée.
\end{definition}

Les indices de référence sont ceux du livre~7 (chapitres 2, 19 et 23)~: le prix moyen pondéré par les volumes, le prix d'arrivée pour
l'implementation shortfall, le cours de clôture. L'article de Madhavan dans l'encyclopédie décrit le VWAP comme un indice de référence courant pour juger le trading et l'exécution, et le manuel de Kissell
(2014) recense les algorithmes que les courtiers ont bâtis sur cet indice. Chaque algorithme est facile à mesurer par rapport à son propre indice de référence, et c'est là le piège~:
l'indice de référence n'est pas ce que le client paie, face au prix auquel l'ordre a été exécuté.

Sur le marché simulé (\texttt{firm.agentmkt} sur \texttt{firm.exchsim}, une journée de 26 intervalles d'une minute dont l'activité suit une courbe en U), chaque
algorithme achète 15\,000 actions, soit environ 5~\% du volume d'une journée ordinaire, sur les mêmes huit jours avec le même flux d'ordres, et chaque jour est également
simulé sans lui. Les prix sont exprimés en ticks, soit un centime pour une action à 100~USD, de sorte qu'un tick vaut un point de base~; les nombres positifs représentent ce que l'acheteur a payé au-dessus de la
référence. Le \Cref{tab:mx:benchmark-algorithms:five} donne l'écart par rapport à l'indice de référence propre à chaque algorithme, de combien ses achats ont déplacé cet indice
(l'indice avec l'algorithme moins l'indice du même jour sans lui), le coût d'impact (exécutions comparées au prix milieu qu'avait le même jour aux mêmes instants,
sans l'algorithme), et le coût par rapport au prix d'arrivée.

\begin{table}[ht]
\centering
\small
\begin{tabular}{@{}p{3.4cm}>{\centering\arraybackslash}p{2.8cm}>{\centering\arraybackslash}p{2.8cm}>{\centering\arraybackslash}p{2.8cm}>{\centering\arraybackslash}p{2.8cm}@{}}
\toprule
algorithme & écart à l'indice (écart type) & indice déplacé & impact (erreur type) & écart au prix d'arrivée (erreur type) \\
\midrule
VWAP & 0.6 (0.2) & 2.0 & 2.7 (0.4) & 3.0 (2.1) \\
TWAP & 0.8 (0.5) & 2.1 & 2.9 (0.2) & 3.3 (2.2) \\
participation 10~\% & 1.5 (0.7) & 3.2 & 4.7 (0.6) & 5.3 (1.7) \\
shortfall, $\kappa T=4$ & 2.9 (5.1) & 0.0 & 2.9 (0.5) & 2.9 (1.8) \\
clôture cible & $-5.9$ (3.6) & 12.6 & 6.8 (1.1) & 7.3 (2.3) \\
\bottomrule
\end{tabular}
\caption{Les cinq algorithmes achetant 15\,000 actions sur les mêmes huit journées simulées (ticks par action)~; l'\omterm{def:mx:benchmark-algorithms:is}{algorithme de clôture cible} envoie 30~\% vers
l'enchère de clôture. Données~:
\texttt{mx\_algos.benchmark\_study}.}\label{tab:mx:benchmark-algorithms:five}
\end{table}

Trois résultats ressortent. L'\omterm{def:mx:benchmark-algorithms:vwap}{algorithme VWAP} manque le VWAP de 0.6 tick avec un écart type de 0.2, le plus serré de tous~; mais ses propres achats ont relevé
le VWAP de la journée de 2.0 ticks, de sorte que le coût qu'il a causé, 2.7 ticks d'impact, dépasse plus de quatre fois son glissement annoncé. L'\omterm{def:mx:benchmark-algorithms:is}{algorithme de clôture cible}
«~bat~» la clôture de 5.9 ticks parce que son \omterm{def:mx:the-limit-order-book:orders}{ordre au marché} à la clôture de 4\,500 actions, croisé contre un carnet de clôture mince, a relevé le cours de clôture de 12.6
ticks~: il a déplacé son propre indice de référence. Par rapport au prix d'arrivée, les mouvements de prix propres à la journée (un écart type d'environ six ticks) noient les
écarts sur huit journées~; seul le coût d'impact à nombres aléatoires communs classe les algorithmes de façon fiable, et là, les \omterm{def:mx:benchmark-algorithms:vwap}{algorithmes de participation} et de clôture
qui négocient le plus vite ou en un seul bloc, coûtent le plus cher.

\section{Prévision du volume intrajournalier}

\begin{definition}[Prévision de volume intrajournalier]\label{def:mx:benchmark-algorithms:forecast}
Une \emph{prévision de volume intrajournalier}\index{prévision de volume intrajournalier} donne, avant l'ouverture ou au fil de la journée, le volume attendu de chaque intervalle
de la journée, ou sa part du total journalier~; le profil de volume intrajournalier (livre~7, chapitre~5) en est la forme la plus simple.
\end{definition}

Le panel simulé compte vingt actions sur 250 journées de 26 intervalles. Chaque action a son propre profil en U, avec une montée en fin de séance vers la clôture~; le volume d'un
intervalle multiplie le niveau de l'action par un facteur de marché (persistant d'une journée à l'autre et autocorrélé à l'intérieur de la journée), un choc propre à la journée, un écart propre à l'action
AR(1) à l'intérieur de la journée et, sur 8~\% des couples action-journée, une nouvelle qui fait bondir le volume depuis un intervalle tiré au hasard puis le fait retomber en environ quatre intervalles. Ces
courbes constituent l'activité de \texttt{firm.agentmkt}~: elles règlent le rythme des ordres de ses \omterm{def:mx:why-there-is-a-spread:informed}{noise traders} et des cotations de ses fournisseurs de liquidité, et le volume
qu'une journée affiche est celui des ordres de ses \omterm{def:mx:why-there-is-a-spread:informed}{noise traders}, que le simulateur tire à l'avance à partir de la graine de la journée. La \Cref{fig:mx:benchmark-algorithms:profile} montre
le profil moyen d'une action, une journée ordinaire et une journée de nouvelle.

\begin{omfigure}
\begin{tikzpicture}
\begin{axis}[width=11cm,height=5.4cm,xlabel={\small intervalle (intervalles d'une minute de la journée simulée)},ylabel={\small part du volume de la journée (\%)},
  tick label style={font=\footnotesize},grid=major,grid style={gray!25},xmin=0.5,xmax=26.5,ymin=0,ymax=16,
  legend columns=3,legend style={font=\footnotesize,at={(0.5,-0.3)},anchor=north,draw=none,fill=none,
  /tikz/every even column/.append style={column sep=8pt}},table/col sep=comma]
\addplot[omDef,very thick] table[x=bin,y=typical]{figdata/microstructure/16-benchmark-algorithms/profile.csv};
\addplot[omThm,thick,mark=*,mark size=1.2pt] table[x=bin,y=normal]{figdata/microstructure/16-benchmark-algorithms/profile.csv};
\addplot[omProp,thick,mark=square*,mark size=1.2pt] table[x=bin,y=news]{figdata/microstructure/16-benchmark-algorithms/profile.csv};
\legend{moyenne sur 250 jours,une journée ordinaire,une journée de nouvelle}
\end{axis}
\end{tikzpicture}

\omcaption{Parts du volume intrajournalier d'une action simulée~: son profil moyen, le U avec une montée vers la clôture, et deux journées isolées. Sur la journée de nouvelle
, le volume bondit à l'intervalle 21, de sorte que la part de chaque intervalle antérieur tombe sous le profil. Données~: \texttt{mx\_algos.flow}.}\label{fig:mx:benchmark-algorithms:profile}
\end{omfigure}

Cinq prévisionnistes sont comparés sur les 190 journées qui suivent une période d'apprentissage de 60 journées, pour chaque action~:
\begin{itemize}
\item la \emph{courbe statique}~: le profil moyen des 60 journées d'apprentissage, jamais mis à jour~;
\item le profil \emph{glissant}~: la moyenne des parts des 20 dernières journées~;
\item \emph{niveau plus AR}~: le profil glissant assorti d'un modèle de la journée, $\log v_k=\log(\bar V p_k)+\ell+e_k$, où le niveau $\ell$ vaut pour toute la journée et
l'écart $e_k$ suit un AR(1)~; après $k$ intervalles, le niveau est estimé par moindres carrés généralisés à partir des écarts observés, et le dernier écart est
laissé s'éteindre (le \cref{lst:mx:benchmark-algorithms:levelar})~;
\item \emph{PCA-ARMA}~: la décomposition du taux de rotation proposée par Białkowski, Darolles et Le Fol (2008), en une composante commune, extraite par analyse en
composantes principales (livre~4, chapitre~22) à travers les actions et prévue par sa moyenne sur les 20 dernières journées, et une composante propre à chaque action prévue
par un modèle ARMA (ils ont aussi utilisé des autorégressions à seuil)~; ici, un AR(1)~;
\item la même, en version \emph{dynamique}~: la composante propre re-prévue après chaque intervalle à partir de la dernière observée.
\end{itemize}

\omcode{code/firm/algos/firm_algos.py}{65}{85}{Le modèle de volume à niveau plus AR~: ses cinq paramètres sont ajustés sur un historique de jours, et après chaque intervalle il réestime le niveau de la journée par moindres carrés généralisés et laisse s'éteindre le dernier écart.}\label{lst:mx:benchmark-algorithms:levelar}

Une erreur de prévision ne compte que par ce qu'elle fait au prix de l'échéancier. Notons $a_k$ les parts de l'échéancier et $b_k$ les parts de volume de la journée, et
$D_k=\sum_{i\le k}(a_i-b_i)$ leur écart cumulé.

\begin{proposition}[Ce que coûte la prévision de volume]\label{prop:mx:benchmark-algorithms:tracking}
Si les VWAP des intervalles varient par accroissements indépendants de variance $s^2$, indépendants des erreurs de l'échéancier, la part du slippage par rapport au VWAP
imputable à l'échéancier, $\sum_k(a_k-b_k)w_k$, a pour variance $s^2\,\E\sum_{k<B}D_k^2$.
\end{proposition}

\begin{proof}
Puisque $\sum_k(a_k-b_k)=0$, la sommation par parties donne $\sum_k(a_k-b_k)w_k=-\sum_{k<B}D_k(w_{k+1}-w_k)$, une somme d'accroissements indépendants pondérés par
$D_k$.
\end{proof}

La tracking error $s\sqrt{\sum D_k^2}$ est celle que rapporte la \cref{fig:mx:benchmark-algorithms:forecast}, pour une action à 2~\% de volatilité quotidienne, sur les 3\,800
jours-action de la période de test et sur les 331 d'entre eux (8.7~\%) dont le volume dépassait 1.5 fois la moyenne de l'action sur les 20 jours précédents. Konishi
(2002) a résolu la version statique du problème, c'est-à-dire le découpage qui minimise cette variance lorsque le volume et le prix sont aléatoires.

\begin{omfigure}
\begin{tikzpicture}
\begin{axis}[width=11cm,height=5.6cm,ybar,bar width=9pt,ylabel={\small tracking error du VWAP (points de base)},
  xtick={0,1,2,3,4},xticklabels={{courbe\\statique},{glissant},{niveau\\+ AR},{PCA-\\ARMA},{PCA-ARMA\\dynamique}},
  xticklabel style={align=center},
  tick label style={font=\footnotesize},ymin=0,ymax=24,grid=major,grid style={gray!25},enlarge x limits=0.14,
  legend columns=2,legend style={font=\footnotesize,at={(0.5,-0.22)},anchor=north,draw=none,fill=none,
  /tikz/every even column/.append style={column sep=8pt}},table/col sep=comma]
\addplot[fill=omDef!70,draw=omDef] table[x=i,y=normal]{figdata/microstructure/16-benchmark-algorithms/forecast.csv};
\addplot[fill=omProp!70,draw=omProp] table[x=i,y=high]{figdata/microstructure/16-benchmark-algorithms/forecast.csv};
\legend{autres jours,jours à fort volume}
\end{axis}
\end{tikzpicture}

\omcaption{La part du slippage par rapport au VWAP imputable à l'échéancier (racine quadratique moyenne de $s\sqrt{\sum D_k^2}$, pour une action à 2~\% de volatilité quotidienne) pour les cinq
prévisionnistes~: 3\,469 autres jours-action et 331 jours à fort volume. Les échéanciers à niveau plus AR et PCA-ARMA dynamique se mettent à jour au cours de la journée. Données~:
\texttt{mx\_algos.forecast\_study}.}\label{fig:mx:benchmark-algorithms:forecast}
\end{omfigure}

Trois enseignements. Prévoir les parts est difficile~: l'erreur totale du meilleur prévisionniste sur les parts de la journée, $\sum_k|a_k-b_k|$, reste de 0.306 (contre 0.324 pour
la courbe statique), car l'essentiel de l'écart d'un intervalle est du bruit que personne ne peut prévoir. Les jours à fort volume coûtent davantage, 17.7 points de base de tracking error pour
la courbe statique contre 13.7 les autres jours, car leur volume supplémentaire arrive par à-coups qui remodèlent la journée. Et la mise à jour aide un peu sur les
jours ordinaires (l'échéancier PCA-ARMA dynamique donne 12.7 contre 13.7) mais peut nuire aux jours qui comptent~: son AR(1) additif, ajusté surtout sur des jours calmes,
prend une poussée liée à l'actualité pour du bruit qui va revenir, et finit à 20.6 points de base, pire que l'absence de mise à jour. Le modèle à niveau plus AR, qui sépare le niveau
d'une journée de ses écarts transitoires, est le seul à faire mieux sur les deux types de jours (13.3 et 17.2). Le profil glissant ne fait pas mieux que
la courbe statique ici, parce que le profil simulé ne dérive pas~; sur des données réelles, il le fait.

\section{VWAP~: statique et dynamique}

\begin{definition}[VWAP dynamique]\label{def:mx:benchmark-algorithms:dvwap}
Un \emph{VWAP dynamique}\index{VWAP dynamique} échéancier recalcule, après chaque intervalle, la quantité restant à exécuter sur les intervalles restants, proportionnellement à une
prévision actualisée de leur volume, au lieu de figer tout l'échéancier avant l'ouverture.
\end{definition}

Le slippage d'un achat par rapport au VWAP de la journée se décompose exactement. Avec $q_k$ actions achetées dans l'intervalle $k$ à un prix moyen $f_k$, le volume de marché $v_k$
et le VWAP $w_k$, $Q=\sum q_k$ et $V=\sum v_k$,
\[
\frac{\sum_kq_kf_k}{Q}-\frac{\sum_kv_kw_k}{V}=\underbrace{\sum_k\frac{q_k}{Q}(f_k-w_k)}_{\text{exécution}}+\underbrace{\sum_k\Bigl(\frac{q_k}{Q}-\frac{v_k}{V}\Bigr)w_k}_{\text{échéancier}}:
\]
la part exécution est ce que les \omterm{def:mx:the-almgren-chriss-framework:parent}{ordres enfants} de chaque intervalle ont payé par rapport au VWAP de cet intervalle (le spread et l'impact), la part échéancier est l'erreur de prévision du
volume, valorisée aux mouvements de la journée. \texttt{firm\_algos.decompose} calcule les deux~; les tests vérifient qu'elles totalisent le slippage au
centime près.

\omcode{code/microstructure/16-benchmark-algorithms/python/mx_algos.py}{283}{290}{Les échéanciers VWAP statique et dynamique transmis à l'agent d'exécution simulé~: le dynamique réestime les intervalles restants à partir du volume des autres intervenants observé jusqu'ici et attribue au prochain intervalle sa part de ce qui reste à acheter.}\label{lst:mx:benchmark-algorithms:schedules}

Humphery-Jenner (2011) a construit un \omterm{def:mx:benchmark-algorithms:dvwap}{VWAP dynamique} qui laisse l'échéancier réagir aux nouvelles au cours de la journée. Ici, les deux échéanciers achètent 15\,000 actions (4.95~\% du volume sur
les jours ordinaires, 2.5~\% sur les jours animés) sur seize jours à fort volume et seize autres jours-action du panel simulé, sur le marché simulé, avec
le flux de la journée identique pour les deux.
La \Cref{fig:mx:benchmark-algorithms:vwap} trace le slippage de chaque jour en fonction de sa part échéancier.

\begin{omfigure}
\begin{tikzpicture}
\begin{axis}[width=8.6cm,height=6.4cm,xlabel={\small part échéancier (ticks par action)},ylabel={\small slippage par rapport au VWAP (ticks)},
  tick label style={font=\footnotesize},grid=major,grid style={gray!25},xmin=-3,xmax=8,ymin=-3,ymax=8,
  legend columns=3,legend style={font=\footnotesize,at={(0.5,-0.26)},anchor=north,draw=none,fill=none,
  /tikz/every even column/.append style={column sep=6pt}},table/col sep=comma]
\addplot[omDef,only marks,mark=*,mark size=1.8pt] table[x=sched,y=total]{figdata/microstructure/16-benchmark-algorithms/vwap_high_static.csv};
\addplot[omProp,only marks,mark=square*,mark size=1.8pt] table[x=sched,y=total]{figdata/microstructure/16-benchmark-algorithms/vwap_high_dynamic.csv};
\addplot[gray,only marks,mark=o,mark size=1.5pt] table[x=sched,y=total]{figdata/microstructure/16-benchmark-algorithms/vwap_normal_static.csv};
\addplot[black,dashed,domain=-3:8] {x};
\legend{statique (fort volume),dynamique (fort volume),statique (autres jours)}
\end{axis}
\end{tikzpicture}

\omcaption{Slippage par rapport au VWAP d'un achat de 15\,000 actions en fonction de sa part échéancier, un point par jour simulé. La droite en pointillés représente le slippage égal à la
part échéancier~; les points se situent à environ 0.7 tick au-dessus, soit le spread et l'impact de la part exécution. Données~: \texttt{mx\_algos.vwap\_study}.}\label{fig:mx:benchmark-algorithms:vwap}
\end{omfigure}

Les jours à fort volume, la part échéancier explique 99.1~\% de la variance du slippage du VWAP statique et 98.6~\% de celle du dynamique~: l'écart
type du slippage est de 2.24 ticks pour l'échéancier statique et 1.75 pour le dynamique, celui de la part échéancier 2.17 et 1.71, tandis que la
part exécution ne varie que de 0.21 et 0.20 autour d'une moyenne de 0.66 et 0.73 ticks (le demi-spread et l'impact que paie chaque \omterm{def:mx:the-almgren-chriss-framework:parent}{ordre enfant}). Les autres jours, le
slippage est plus faible (écarts types 0.48 et 0.41) et la part échéancier en explique moins, 73.9~\% et 68.5~\%. La moindre
dispersion de l'échéancier dynamique sur les jours à fort volume n'est pas un gain fiable sur seize jours~: il a fait mieux sur sept d'entre eux, l'écart apparié des slippages au carré
est de 1.5, avec une erreur type de 1.8, et l'essentiel vient d'un seul jour d'actualité (7.5 ticks pour l'échéancier statique, 5.4 pour le dynamique). Les
3\,800 jours du panel, où la tracking error passe de 17.7 à 17.2 points de base, mesurent le même effet avec précision~: réel, et petit.

\section{Participation et rétroaction}

Un \omterm{def:mx:benchmark-algorithms:vwap}{algorithme de participation} négocie à chaque intervalle une part $r$ du volume de l'intervalle précédent. Le volume qu'il observe comprend ses propres transactions, et cela le conduit à
courir après lui-même~: s'il négocie $q$ contre les $o$ des autres négociants à chaque intervalle, alors $q=r(o+q)$, donc il négocie $q/o=r/(1-r)$ du volume que le marché
aurait eu sans lui, 25~\% pour $r=20~\%$. Deux algorithmes de ce type achetant la même action se comptent mutuellement~: $q=r(o+2q)$, donc chacun négocie
$r/(1-2r)$, 33~\%, et le volume total croît de $1/(1-2r)$. Et si les autres négociants réagissent à ses achats (ici, les fondamentalistes vendent dans le
cours qu'il fait monter), leur volume supplémentaire alimente la même boucle.

\omcode{code/microstructure/16-benchmark-algorithms/python/mx_algos.py}{328}{335}{Un échéancier de participation pour l'agent simulé~: à chaque intervalle, une part du volume total du dernier intervalle, ses propres transactions comprises, jusqu'à ce que l'ordre soit exécuté.}\label{lst:mx:benchmark-algorithms:pov}

Le \Cref{tab:mx:benchmark-algorithms:pov} simule un et deux acheteurs de ce type, pour 15\,000 actions à $r=20~\%$, sur une journée normale et sur une journée peu active, avec 40~\% de
l'activité normale, huit jours chacun, contre les mêmes jours sans eux.

\begin{table}[ht]
\centering
\small
\setlength{\tabcolsep}{4pt}%
\begin{tabular}{@{}llrrr>{\centering\arraybackslash}p{2.6cm}>{\centering\arraybackslash}p{2.2cm}@{}}
\toprule
jour & acheteurs & minutes & participation & rythme & volume supplémentaire des autres (e.t.) & coût d'impact (e.t.) \\
\midrule
normal & un & 5.2 & 17.0~\% & 21.1~\% & 0.12 (0.10) & 6.4 (0.7) \\
normal & deux & 4.6 & 15.9~\% & 24.6~\% & 0.12 (0.04) & 10.0 (0.9) \\
peu actif & un & 11.6 & 19.7~\% & 24.9~\% & 0.04 (0.06) & 5.4 (0.8) \\
peu actif & deux & 9.5 & 19.2~\% & 29.4~\% & $-0.05$ (0.04) & 5.7 (0.6) \\
\bottomrule
\end{tabular}
\caption{\omterm{def:mx:benchmark-algorithms:vwap}{Algorithmes de participation} à 20~\% sur des jours simulés normaux et peu actifs~: temps nécessaire pour acheter 15\,000 actions, participation (volume propre rapporté au volume total),
rythme (volume propre rapporté au volume que la même journée aurait eu sans les algorithmes), volume supplémentaire des autres par action achetée, et coût d'impact en ticks.
Pour deux acheteurs, les chiffres sont ceux du premier. Données~: \texttt{mx\_algos.pov\_study}.}\label{tab:mx:benchmark-algorithms:pov}
\end{table}

La participation reste proche de sa cible~; le rythme, lui, non. Sur la journée normale, l'ordre dure cinq minutes, trop peu pour que la boucle se mette en place (la
dernière tranche est aussi partielle), et le rythme vaut 21.1~\%~; sur la journée peu active, il dure presque douze minutes et atteint 24.9~\%, soit le $r/(1-r)$ de l'argument. Avec un
second acheteur, le rythme monte à 24.6~\% et 29.4~\%, vers $r/(1-2r)$, et sur la journée normale le coût d'impact passe de 6.4 à 10.0 ticks~: chaque algorithme paie les achats de l'autre et
s'exécute plus vite à cause de cela. La réaction des autres traders ajoute peu de volume ici (0.12 action par action achetée ou moins,
à peine mesurable). Le coût par action plus bas de la journée peu active tient au marché lui-même~: ses fondamentalistes ramènent le cours à un rythme qui ne
s'amenuise pas avec la journée, et un ordre plus lent leur laisse davantage de temps. Une vraie journée peu active coûte en général davantage~; ce qui se transpose, c'est que l'horloge de l'algorithme est
le volume, et que, sur une journée peu active, ses propres transactions pèsent davantage dans cette horloge.

\section{Algorithmes implementation shortfall et de clôture}

L'\omterm{def:mx:benchmark-algorithms:is}{algorithme d'implementation shortfall} du \cref{tab:mx:benchmark-algorithms:five} suit la trajectoire du chapitre~14 avec $\kappa T=4$~: il avance ses achats,
achetant 87~\% de l'ordre dans la première moitié de la journée. Son indice de référence est le prix d'arrivée, si bien qu'il est le seul algorithme mesuré contre le prix
que l'ordre a effectivement obtenu~; son coût face au prix d'arrivée, 2.9 ticks, a le plus grand écart parmi les chiffres mesurés contre un indice de référence, seulement parce que les
indices de référence des autres évoluent avec la journée. L'\omterm{def:mx:benchmark-algorithms:is}{algorithme de clôture cible} (target close) réserve 30~\% pour l'enchère de clôture (un \omterm{def:mx:the-limit-order-book:orders}{ordre au marché} à la clôture envoyé 30 secondes avant la
clôture) et exécute le reste sur les neuf derniers intervalles, au prorata du profil. Son suivi serré de la clôture se paie par le coût d'impact le plus élevé
(6.8 ticks), et son gain apparent face à la clôture n'est que la clôture qu'il a lui-même déplacée.

Ce que le tableau ne montre pas, c'est le choix entre eux, qui revient au gérant de portefeuille~: un indice de référence VWAP convient à un ordre dont l'\omterm{def:mx:the-almgren-chriss-framework:urgency}{urgence} est
faible et dont le trader veut ressembler à la moyenne de la journée~; l'implementation shortfall convient à un ordre dont l'alpha se dégrade~; la clôture convient aux fonds
valorisés à la clôture. Le chapitre~19 compare les algorithmes entre eux par l'analyse des coûts de transaction (TCA), et le chapitre~28 en construit un qui exécute chaque
tranche par le placement et l'acheminement des chapitres 17 et 18.

\section{Tutoriel~: prévisions de volume et algorithmes de référence}\label{tut:mx:benchmark-algorithms:tut}

\begin{tutorial}
\textbf{Objectif.} Prévoir le volume intrajournalier de trois façons, faire tourner les cinq algorithmes sur le marché simulé, et décomposer le slippage par rapport au VWAP en sa
composante calendrier et sa composante exécution.
\textbf{État final~:} les \cref{fig:mx:benchmark-algorithms:profile,fig:mx:benchmark-algorithms:forecast,fig:mx:benchmark-algorithms:vwap} et les deux tableaux.
\begin{tutsteps}
\item \textbf{Panneau.} \texttt{mx\_algos.panel()} et \texttt{flow()}~: les courbes d'activité et le volume qu'affiche le simulateur.
\item \textbf{Prévisions.} \texttt{firm\_algos.Profile}, \texttt{LevelAR.fit(history)}, \texttt{PCAARMA(window).model(j)}~; \texttt{forecast\_study()}.
\item \textbf{Agent.} \texttt{Slicer(schedule)} sur \texttt{firm.agentmkt}, dans sa \texttt{session}~; \texttt{run}, \texttt{day\_stats}, \texttt{fills}.
\item \textbf{VWAP.} \texttt{vwap\_schedules(j, d)}, \texttt{vwap\_study()} avec \texttt{firm\_algos.decompose}.
\item \textbf{Participation et le reste.} \texttt{pov\_study()}, \texttt{benchmark\_study()}~; tracer les figures avec \texttt{fig\_algos.py}.
\end{tutsteps}
\textbf{Pour aller plus loin.} Faire monter la volatilité du prix avec le volume (les nouvelles font bouger les cours), ce qui rend la composante calendrier plus
grande précisément les jours où la prévision est la moins bonne~; ajouter un second facteur ACP~; doter l'\omterm{def:mx:benchmark-algorithms:vwap}{algorithme de participation} d'un cap et d'un floor sur son taux de participation.
\end{tutorial}

\section{Construction~: algorithmes de référence}\label{bld:mx:benchmark-algorithms:algos}

\begin{build}
\bfield{Rôle} Les calendriers de référence et les modèles de prévision de volume qu'appellent l'\omterm{def:mx:benchmark-algorithms:algo}{algorithme d'exécution} du chapitre~28, l'analyse des coûts de transaction du
chapitre~19 et le desk du chapitre~20.
\bfield{Interface} \texttt{profile(history)}, \texttt{Profile(prof, adv)}, \texttt{LevelAR(prof, adv, rho, s\_level, s\_dev)} et
\texttt{LevelAR.fit(history)}, \texttt{PCAARMA(window, r)} et \texttt{.model(j)}, chacun avec \texttt{remaining(observed)}~;
\texttt{vwap(Q, model, day, dynamic)}, \texttt{twap(Q, bins)}, \texttt{pov(rate, others, own\_counted)}, \texttt{shortfall(Q, bins, kT)} (sur
\texttt{firm.acexec}), \texttt{close(Q, model, moc\_share, start)}, \texttt{decompose(q, fills, v, vwaps)}, \texttt{tracking(q, v)}.
\bfield{Règles} Volumes par intervalle, jours $\times$ intervalles pour les historiques~; calendriers en actions par intervalle, dont la somme égale l'ordre~; un calendrier dynamique ne voit que les intervalles
déjà terminés~; slippage par action, positif pour un acheteur ayant payé davantage.
\bfield{Tests d'acceptation} \texttt{code/firm/algos/tests/}~: le calendrier dynamique d'un modèle statique égale son calendrier statique~; un intervalle à fort volume fait pencher le suivant, et une
journée chargée conforme au profil ne le fait pas~; l'ajustement level-plus-AR retrouve ses paramètres~; PCA-ARMA retrouve une forme commune~; l'état stationnaire de la participation~; la
décomposition est additive.
\bfield{Pour aller plus loin} Prévisions de volume avec un régime pour les jours d'actualité~; une bande de participation (chapitre~28)~; une prévision du déséquilibre à l'enchère de clôture.
\end{build}

\begin{omsources}
\item H.~Konishi, «~Optimal slice of a VWAP trade~», \emph{Journal of Financial Markets} 5(2), 2002.
\item J.~Białkowski, S.~Darolles et G.~Le~Fol, «~Improving VWAP strategies~: a dynamic volume approach~», \emph{Journal of Banking and Finance}
32(9), 2008.
\item A.~Madhavan, «~Volume-weighted average price (VWAP)~», \emph{Encyclopedia of Quantitative Finance}, Wiley, 2010.
\item M.~Humphery-Jenner, «~Optimal VWAP trading under noisy conditions~», \emph{Journal of Banking and Finance} 35(9), 2011.
\item R.~Kissell, \emph{The Science of Algorithmic Trading and Portfolio Management}, Academic Press, 2014.
\end{omsources}

\section{Exercices}

\begin{exercise}[$\star$]\label{exo:mx:benchmark-algorithms:1}
Le profil moyen de l'action simulée attribue 5.50~\% du volume de la journée au premier intervalle et 9.38~\% au dernier. Combien d'actions d'un ordre VWAP statique de
15\,000 actions vont à chacun~?
\end{exercise}

\begin{exercise}[$\star$]\label{exo:mx:benchmark-algorithms:2}
On demande à un \omterm{def:mx:benchmark-algorithms:vwap}{algorithme de participation} de viser 10~\% du volume, mais il ne suit que le volume des autres traders et n'en exécute que 10~\%. Quelle part de l'ensemble du
volume atteint-il, et quel taux appliqué au volume des autres donnerait 10~\%~?
\end{exercise}

\begin{exercise}[$\star$]\label{exo:mx:benchmark-algorithms:3}
Un acheteur exécute 2 actions dans chacun de deux intervalles, en payant 10.10 et 11.00~; le marché a échangé 1 et 3 actions à ces intervalles, à des VWAP de 10.00 et 11.00. Décomposer le
slippage par rapport au VWAP de la journée en ses parts d'exécution et de calendrier.
\end{exercise}

\begin{exercise}[$\star\star$]\label{exo:mx:benchmark-algorithms:4}
Pourquoi la composante calendrier explique-t-elle 99~\% de la variance du slippage par rapport au VWAP les jours de fort volume, mais seulement environ 70~\% les autres jours~?
\end{exercise}

\begin{exercise}[$\star\star$]\label{exo:mx:benchmark-algorithms:5}
L'\omterm{def:mx:benchmark-algorithms:is}{algorithme de clôture cible} a battu la clôture de 5.9 ticks et affiche le coût le plus élevé face au prix d'arrivée. Concilier les deux résultats.
\end{exercise}

\begin{exercise}[$\star\star$]\label{exo:mx:benchmark-algorithms:6}
Trois \omterm{def:mx:benchmark-algorithms:vwap}{algorithmes de participation} à $r=20\%$ achètent la même action, chacun en suivant le volume total. À quel rythme chacun exécute-t-il, par rapport au volume
que le marché aurait eu sans eux~? Que se passe-t-il à $r=25~\%$~?
\end{exercise}

\begin{exercise}[$\star\star\star$]\label{exo:mx:benchmark-algorithms:7}
\emph{Programmation.} Relancer l'étude de prévision avec des nouvelles sur 20~\% des jours-action au lieu de 8~\%. Combien de jours-action sont à fort volume, et quelles sont les
tracking errors de la courbe statique et du calendrier level-plus-AR sur ces jours~?
\end{exercise}

\begin{exercise}[$\star\star\star$]\label{exo:mx:benchmark-algorithms:8}
\emph{Trouver la faille.} «~Notre \omterm{def:mx:benchmark-algorithms:vwap}{algorithme VWAP} a fini à moins de 0.6 ticks du VWAP en moyenne, il a donc coûté 0.6 ticks à notre client.~»
\end{exercise}

\section{Problème~: le VWAP un jour qui n'était pas moyen}

\begin{problem}[{Problème du week-end --- le VWAP un jour qui n'était pas moyen}]\label{pb:mx:benchmark-algorithms:1}
Un client demande pourquoi ses ordres VWAP s'écartent le plus de leur indice de référence les jours où l'action est la plus active, et si un algorithme dynamique améliorerait les choses.

\medskip\noindent\textbf{Partie I --- Les algorithmes.}
\begin{enumerate}
\item Définir les \omterm{def:mx:benchmark-algorithms:vwap}{algorithmes VWAP}, TWAP, de participation, d'implementation shortfall et de clôture cible, ainsi que leurs indices de référence.
\item Quel effet les achats de l'\omterm{def:mx:benchmark-algorithms:vwap}{algorithme VWAP} ont-ils eu sur le VWAP de la journée sur le marché simulé, et quel a été son coût d'impact~?
\item Pourquoi l'\omterm{def:mx:benchmark-algorithms:is}{algorithme de clôture cible} semble-t-il battre son indice de référence~?
\item Quelle mesure classe de façon fiable les cinq algorithmes sur huit jours, et pourquoi~?
\end{enumerate}

\medskip\noindent\textbf{Partie II --- La prévision du volume.}
\begin{enumerate}[resume]
\item Décrire le panel simulé et ce qui fait une journée à fort volume.
\item Énoncer le modèle niveau plus AR et sa mise à jour au cours de la journée.
\item Démontrer que la variance de la partie calendrier vaut $s^2\E\sum D_k^2$.
\item Donner les erreurs de suivi de la courbe statique et du calendrier niveau plus AR les jours à fort volume et les autres jours.
\item Pourquoi le calendrier dynamique PCA-ARMA fait-il moins bien que le calendrier statique les jours à fort volume~?
\end{enumerate}

\medskip\noindent\textbf{Partie III --- Le VWAP sur le marché simulé.}
\begin{enumerate}[resume]
\item Écrire la décomposition exacte du slippage par rapport au VWAP en partie exécution et partie calendrier.
\item Comment les jours à fort volume ont-ils été choisis, et pourquoi les deux calendriers voient-ils le même flux d'ordres~?
\item Donner les écarts types du slippage, de sa partie calendrier et de sa partie exécution pour les deux calendriers les jours à fort volume.
\item Quelle est la moyenne de la partie exécution, et d'où vient-elle~?
\item \emph{Énoncer le résultat nommé}~: quelle part du slippage par rapport au VWAP s'explique par l'erreur de prévision du volume,
pour le VWAP statique et le \omterm{def:mx:benchmark-algorithms:dvwap}{VWAP dynamique}, les jours à fort volume du simulateur~?
\item Le gain du calendrier dynamique est-il fiable~? Donner les éléments de preuve.
\end{enumerate}

\medskip\noindent\textbf{Partie IV --- La participation et la réponse.}
\begin{enumerate}[resume]
\item Dériver le rythme $r/(1-r)$ d'un \omterm{def:mx:benchmark-algorithms:vwap}{algorithme de participation} qui compte ses propres transactions.
\item Que fait un second acheteur au rythme et au coût le jour normal~?
\item Pourquoi le rythme était-il plus élevé le jour peu liquide~?
\item Que diriez-vous au client au sujet du \omterm{def:mx:benchmark-algorithms:dvwap}{VWAP dynamique}~?
\item En une phrase~: de quoi le slippage d'un algorithme de référence est-il une mesure~?
\end{enumerate}
\end{problem}

\section{Questions d'entretien}

\begin{interviewq}[$\star$ \iqroles{trader}]\label{iq:mx:benchmark-algorithms:1}
Un client veut acheter 5~\% du volume de la journée sur une action liquide, sans opinion particulière. Quel algorithme choisiriez-vous, et pourquoi~?
\end{interviewq}

\begin{interviewq}[$\star\star$ \iqroles{chercheur}]\label{iq:mx:benchmark-algorithms:2}
Comment prévoiriez-vous la courbe de volume intrajournalier d'une action, et comment jugeriez-vous la prévision~?
\end{interviewq}

\begin{interviewq}[$\star\star$ \iqroles{chercheur}]\label{iq:mx:benchmark-algorithms:3}
Pourquoi le slippage d'un \omterm{def:mx:benchmark-algorithms:vwap}{algorithme VWAP} ne dépend-il de la prévision de volume qu'à travers les erreurs cumulées~?
\end{interviewq}

\begin{interviewq}[$\star\star$ \iqroles{trader}]\label{iq:mx:benchmark-algorithms:4}
Votre \omterm{def:mx:benchmark-algorithms:vwap}{algorithme de participation} est à 20~\% sur une action où un autre gros acheteur utilise aussi un algorithme de ce type. Que se passe-t-il, et que feriez-vous~?
\end{interviewq}

\begin{interviewq}[$\star\star$ \iqroles{développeur}]\label{iq:mx:benchmark-algorithms:5}
Concevez l'interface entre un calendrier et la couche des \omterm{def:mx:the-almgren-chriss-framework:parent}{ordres enfants} d'un \omterm{def:mx:benchmark-algorithms:algo}{algorithme d'exécution}.
\end{interviewq}

\begin{interviewq}[$\star\star\star$ \iqroles{banque}]\label{iq:mx:benchmark-algorithms:6}
Un client se plaint que votre \omterm{def:mx:benchmark-algorithms:is}{algorithme de clôture cible} ait battu la clôture, mais que la performance du fonds ait souffert. Qu'a-t-il pu se passer~?
\end{interviewq}
